% ============================================
% Supplementary Note 2: Shared Assembly History Promotes Dominance
% ============================================
\subsection*{Supplementary Note 2: Shared assembly history promotes Dominance}
\label{sec:suppl:note2}

To test whether assembly history contributes to community-level selection~\citep{Gilpin1994,Debray2022,Tikhonov2016,DiazColunga2022}, we compared coalescence outcomes to a ``direct assembly'' comparison. \revsecond{This comparison asks whether the Dominance signature reflects excess origin-correlated persistence generated by prior assembly, rather than only assigning a categorical label to the observed outcome.} In coalescence, two parental communities are first assembled separately from the species pool and then mixed. In direct assembly, all species from both pools are mixed simultaneously without prior assembly history.

In simulations, species compete during assembly and some go extinct. Those that survive to coexistence tend to have weaker mutual interactions on average. We quantified this assembly effect by measuring the average pairwise interaction coefficient $\langle \alpha_{ij} \rangle$ among species at different stages (Supplementary Fig.~32) and by visualizing pre- and post-assembly interaction submatrices (Supplementary Fig.~3). Across all tested interaction-strength parameter values ($\mu = 0$ to $1.2$), post-assembly communities exhibited substantially reduced average pairwise interaction coefficients compared to the pre-assembly pool; at the representative value $\mu = 0.6$, this reduction was significant across 600 paired pre- versus post-assembly summaries (paired $t$-test, $t_{599} = 29.26$, $p = 1.54 \times 10^{-117}$; Fig.~2C). This reduction becomes more pronounced at higher $\mu$, where stronger competition drives more species to extinction, leaving only those with weaker mutual interactions. The assembly process thus creates internal coherence within each parental community: species that survived together share compatible (less competitive) interactions.

We hypothesized that this shared history would increase Dominance probability compared to direct assembly. Simulations supported this prediction (Extended Data Fig.~4): across interaction-strength parameter values $\mu = 0.2$ to $1.0$, coalescence consistently produced higher Dominance fractions than direct assembly. At $\mu = 0.2$, coalescence yielded 37\% Dominance versus 12\% for direct assembly. At $\mu = 0.6$, the values were 59\% versus 20\%. At $\mu = 1.0$, they were 65\% versus 22\%. Correspondingly, direct assembly showed higher Restructuring rates, as expected when species from different pools have not been pre-filtered for compatibility.

Experimental results corroborate these findings (Supplementary Fig.~23). Under Base and Nutr$+$ conditions, coalescence showed elevated Dominance rates compared to direct assembly, consistent with simulation predictions. The Nutr$-$ condition showed a weaker coalescence-versus-direct-assembly contrast, consistent with weaker invasion resistance in nutrient-limited environments (see Supplementary Note~4). A separate pair-level analysis of parental-affiliation-correlated species fates is presented in Supplementary Note~7.
